\section{Gate Simulation and Statistical Power}
\label{app:sim}

We chose rule P before any main \capevolve{} seed, from a budget-exact simulation on the real promotion and gate code. We then re-ran it at each suite's observed difficulty: the seed artifact's mean reward and trial-level variance on the search split (Suite A: 40 validation tasks, 600 rollouts; Suite B: 49 optimization tasks, 550 rollouts; Suite C: 20 tasks), one screening trial per task, about ten candidates per search and 22{,}400 searches per cell. The pre-declared criterion was to maximise the expected true new-task gain of the final artifact subject to at most 0.15 probability of a false promotion in an all-null search. A search screens 6.7 to 9.4 candidates on average. Table~\ref{tab:sim} gives rule P and the one-stage gate. The false-promotion probability is per search, about one percent per candidate. The gain column is the expected true new-task gain of the final artifact when 60\% of proposed edits are null. The one-stage gate has the larger expected gain and the higher power, and it falsely promotes in about half of all-null searches; rule P gives up power to avoid that. Five confirmation trials at $z_{\mathrm{conf}}=1.645$ would cost $2\cdot 40\cdot 5=400$ of the 600 rollouts per confirmation on Suite A instead of 240, so fewer candidates are screened and the expected gain halves; on Suite B the 490 rollouts exceed the 452 left after the baseline and one screen, so that setting could never promote.

\begin{table}[h]
\caption{Rule P ($n_{\mathrm{conf}}=3$, $z_{\mathrm{conf}}=1.282$) and the one-stage $k=1$ gate, per suite, from simulation at observed difficulty. False prom.\ is the probability of at least one false promotion in an all-null search of about ten candidates. The four power columns give the probability of promoting one candidate with the stated true gain. Gain is the expected true new-task gain of the final artifact.}
\label{tab:sim}
\centering\small
% GENERATED by scripts/derive_numbers.py. Do not edit by hand.
\begin{tabular}{@{}lrrrrr@{}}
\toprule
Rule & False prom. & $+0.04$ & $+0.06$ & $+0.08$ & $+0.10$ \\
\midrule
P, $n_c{=}3$, $z{=}2.0$, $\delta_{\min}{=}0.02$ & 0.014 & 0.03 & 0.08 & 0.15 & 0.26 \\
P, $n_c{=}3$, $z{=}1.645$ & 0.030 & 0.06 & 0.12 & 0.21 & 0.33 \\
P, $n_c{=}3$, $z{=}1.282$ (used) & 0.055 & 0.09 & 0.17 & 0.27 & 0.41 \\
P, $n_c{=}5$, $z{=}1.645$ & 0.019 & 0.08 & 0.16 & 0.29 & 0.42 \\
one-stage, $k{=}1$ & 0.565 & 0.29 & 0.37 & 0.45 & 0.54 \\
one-stage, $k{=}1.645$ & 0.297 & 0.12 & 0.17 & 0.23 & 0.29 \\
\bottomrule
\end{tabular}

\end{table}

Seed-artifact variance replaces the earlier assumptions (mean 0.55, within-task variance 0.149). The observed within-task variance is 0.088, 0.148 and 0.124 on Suites A, B and C, and the standard deviation of the true pass rate across tasks is 0.40, 0.31 and 0.37. With three search seeds, 80\% power, $\alpha=0.05$ and a task-level effect spread of 0.15, the minimum detectable effect for contrasts between the three frozen artifacts is \mdeACondOne{} on Suite A (\mdeATasks{} tasks, one trial), \mdeBCondOne{}, \mdeBCondTwo{} and \mdeBCondThree{} on Suite B (40 tasks, one to three trials) and \mdeCCondOne{}, \mdeCCondTwo{} and \mdeCCondThree{} on Suite C (20 tasks). For claims about the optimizer, which include seed variance and use $t$ with two degrees of freedom, it is \mdeAPipeOne{} on A, \mdeBPipeOne{} to \mdeBPipeThree{} on B and \mdeCPipeOne{} to \mdeCPipeThree{} on C. With Holm correction over the primary family the first-step values rise to \mdeAHolmOne{} on A, \mdeBHolmOne{} on B and \mdeCHolmOne{} on C with one trial. The task-level effect spread, seed-by-task interaction and seed variance cannot be estimated from seed-artifact data and remain assumptions.

The shared standard-error helper adds a within-task term to a variance of task means that already contains it. In simulation (mean reward 0.55, dispersion 1.5, 40 to 280 tasks) its variance is 1.00 times the correct value at one trial per task, 1.33 times at three trials, 1.24 at five and 1.13 at ten for the new-task estimand, and 1.6 to 8.6 times too large for the fixed-task estimand. Its coverage of a new-task target is 0.94 to 0.97. At one trial it equals the between-task term alone.

\section{Scorer Audit}
\label{app:scorer}

The default SpreadsheetBench reward is \texttt{hard\_no\_recalc}: cached values as saved. Correct formula outputs without cached values passed 0 of 30 checks under it and 25 of 30 (LibreOffice 7.1.8.1) to 29 or 30 of 30 (LibreOffice 25.8.7.3) under recalculation. The recalculation path recalculates the candidate output only and uses gold as shipped. With LibreOffice 7.1.8.1 the official gold file used as output fails on 32 of 400 verified tasks, and with 25.8.7.3 on 19. Four verified tasks (130-9, 283-32, 45944, 49300) are unscorable on every path because of malformed answer-position strings, all four in the SkillOpt test split. Scoring was deterministic: 333 records were scored twice with fresh profiles and none differed. The gold workbooks sat in the directory bind-mounted read-write into the agent container; the campaign runs agents on a gold-free copy and scores on the pristine tree.

\section{Historical Results Under Both Scorers}
\label{app:hist}

\begin{table}[h]
\caption{Historical SpreadsheetBench runs (SkillOpt split, 280 test tasks, three optimization runs per row, Opus 5 optimizer) recomputed from run artifacts. Test-280 was development-exposed, so these are fit-quality figures, not held-out results. Paired means; the as-saved delta is the anti-conservative pooled-observation figure of the earlier draft.}
\label{tab:hist}
\centering\small
\begin{tabular}{@{}lrrrr@{}}
\toprule
Target & Baseline as saved & Gain as saved & Baseline recalc. & Gain recalc. \\
\midrule
Gemma-4-31B & \histGemmaBaseSaved{} & +\histGemmaDeltaSaved{} & \histGemmaBaseRecalc{} & +\histGemmaDeltaRecalc{} \\
GPT-5.5 & \histGptBaseSaved{} & +\histGptDeltaSaved{} & \histGptBaseRecalc{} & +\histGptDeltaRecalc{} \\
\bottomrule
\end{tabular}
\end{table}

Disputed and withdrawn claims: the earlier draft quoted 76.5 for Gemma while a blog post quoted 76.8 for the same data (a denominator definition: the post counts the champion on all champion-scored tasks, including four whose seed rollout failed for infrastructure reasons). A GPT-5.5 run with a champion test score of 89.8 against a baseline of 28.9 is excluded: it is a 225-task subset in which 55 champion rollouts failed on a budget error, and the seed scores 38.6 on all 280. Full lists are in \texttt{audit/recomputed\_numbers.csv}.

\section{Infrastructure Incidents}
\label{app:incidents}

Times are UTC on 2026-10-11 unless stated. Each incident lists its effect on the data.
\begin{itemize}[leftmargin=*,nosep]
\item[\emph{EvoSkill priority starvation.}] The gateway's low-priority tag list still contained EvoSkill after its promotion to Tier 1, so under reserved capacity its calls never got a slot between 13:05 and 18:40. Those rollouts show as 1800-second timeouts with zero model calls. They are invalid and kept as \texttt{INVALID} directories; EvoSkill restarted from scratch at 18:25 to 18:55 and is about 5.5 hours behind the other arms.
\item[\emph{Suite B relaunch.}] The first Suite B runs (09:40) used the wrong judge, a pre-merge worktree without rule P, and untagged candidate-tool errors. They were stopped at 12:38, labelled superseded and not used. A second correction restored the agent-optimize seed capability and archived two sessions that ran the archived legacy engine.
\item[\emph{Content-filter rename for ACE.}] Through the gateway, Claude Sonnet 5.5 returns an empty content-filter response when a prompt asks for a chain-of-thought or reasoning field in the output; prose mentions pass. ACE's prompts requested such a field. Its output field was renamed to ``analysis'' with the diff and affected-call count recorded. No other arm was affected.
\item[\emph{HarnessLens retry policy.}] Upstream aborts a whole baseline stage on one worker error. One of six trial rollouts ended with \texttt{opencode\_exit\_1} after a streamed gpt-oss response ended with only reasoning content and no finish reason. HarnessLens is the only arm that streams the target (1 of 98 streamed calls lacked usage). The campaign policy retries a worker error once inside the rollout budget, scores a second error as zero for HarnessLens's internal selection only, and logs every occurrence.
\item[\emph{Gateway call failures.}] At 13:11 we added a transparent retry (three retries, exponential backoff) for connection resets, read timeouts and 502 to 504 responses, and raised the upstream timeout from 900 to 1800 s. Before that, 1 to 2.5\% of Gemma calls failed, and the shared SpreadsheetBench adapter scored such failures as zero for competitor arms while \capevolve{} dropped them. Pre-13:11 zeros remain in competitor histories and are labelled.
\item[\emph{Gold and split-file readability.}] On Suite A the shared proposer user could read split-id files in world-readable checkouts until 10:00; gold and credentials were already denied. A Suite B agent-optimize transcript scan found no read of task or split files, and the arm was moved into a mount-namespace shim at 15:00. Meta-Harness first launched with a gold-readable proposer, stopped before any proposer iteration, and was relaunched container-isolated at 09:38.
\item[\emph{Failed reproductions.}] VERSE needs SWE-rebench or Terminal-Bench docker sessions; CoEvoSkills needs a Harbor per-task loop; HarnessLens on Suite A evolves only OpenCode, Codex and Pi runtimes. Blocker files are in the campaign scratch store.
\end{itemize}

\paragraph{Rollouts scored zero after an infrastructure error.}
Before the 3-hour hang rule, ACE and MIPROv2 first ran with a 40-minute watchdog that scored a timeout as reward zero.
The coordinator rejected it as asymmetric before it produced results: the six 40-minute runs (three seeds per arm) were killed after one counted rollout each, and their event logs hold no watchdog zero.
Zeros that did reach the search come from other adapter paths.
An adapter returns reward 0 with an error string when an LLM call times out at the client or the sandbox returns HTTP 500.
Each such rollout counts against the budget, so the search saw a score and not a missing cell.
We did not check whether the ACE reflector or the MIPROv2 bootstrap showed the error text.

Adapters label these errors differently, so we classify every error row with one rule and keep the adapter label beside it (Table~\ref{tab:zeros}).
Search-time feedback used each method's native handling, and only the reporting is unified.
The rule changes Trace2Skill, WikiSkill, Meta-Harness and GEPA, and leaves SkillOpt as labelled.
Trace2Skill and WikiSkill label all their errors as candidate faults, yet the rule moves 26 and 8 of them to infrastructure: the stored Trace2Skill strings are 8 client timeouts and 18 sandbox 500 errors, and the WikiSkill errors made no target call.
Meta-Harness keeps 27 of its 30 infrastructure labels.
SkillOpt has 27 infrastructure errors and 120 candidate errors, and the 120 are rows where the task template broke its placeholder contract, with no target call.
GEPA's 7 infrastructure labels become candidate errors because they have no string and an unknown call count, and the same happens to 3 Meta-Harness rows and a few Suite B and C rows.
These 18 rows are listed as unverifiable, and raw-log strings would let us reclassify them.
GEPA persists outputs only for its best program, so its counts are incomplete.
The one live timeout zero after the 13:11 gateway fix is MIPROv2 seed 0 at 13:20.
Every other timeout zero before 13:11 enters the sensitivity analysis that treats such rollouts as missing.

\begin{table*}[t]
\centering\scriptsize
\setlength{\tabcolsep}{4pt}
\caption{Error rows per arm and seed under each adapter's own label and under one uniform rule (\texttt{error\_label\_parity.csv}). The rule calls an error infrastructure if it has a transport, timeout or sandbox-HTTP string, or made no target call and kept no string. Every other error is a candidate error, except rows an adapter labelled infrastructure that have no recoverable string and some target calls, which stay unverifiable. Changed counts rows whose class differs from the adapter label. Search used each method's native handling of these errors, and only the reporting is unified. Live runs only.}
\label{tab:zeros}
\begin{tabular}{@{}llrrrrrrrr@{}}
\toprule
Suite & Arm & Seed & Errors & Adapter infra & Adapter candidate & Uniform infra & Uniform candidate & Unverifiable & Changed \\
\midrule
A & ACE & 1 & 1 & 1 & 0 & 1 & 0 & 0 & 0 \\
A & ACE & 2 & 1 & 1 & 0 & 1 & 0 & 0 & 0 \\
A & CapEvolve AO & 0 & 1 & 1 & 0 & 1 & 0 & 0 & 0 \\
A & CapEvolve AO & 1 & 2 & 2 & 0 & 2 & 0 & 0 & 0 \\
A & CapEvolve AO & 2 & 2 & 2 & 0 & 2 & 0 & 0 & 0 \\
A & EvoSkill & 0 & 2 & 2 & 0 & 2 & 0 & 0 & 0 \\
A & EvoSkill & 1 & 2 & 2 & 0 & 2 & 0 & 0 & 0 \\
A & EvoSkill & 2 & 1 & 1 & 0 & 1 & 0 & 0 & 0 \\
A & GEPA & 0 & 1 & 1 & 0 & 1 & 0 & 0 & 0 \\
A & GEPA & 1 & 4 & 4 & 0 & 4 & 0 & 0 & 0 \\
A & GEPA & 2 & 2 & 2 & 0 & 2 & 0 & 0 & 0 \\
A & CapEvolve HC & 0 & 2 & 2 & 0 & 2 & 0 & 0 & 0 \\
A & CapEvolve HC & 1 & 1 & 1 & 0 & 1 & 0 & 0 & 0 \\
A & CapEvolve HC & 2 & 3 & 3 & 0 & 3 & 0 & 0 & 0 \\
A & Meta-Harness & 0 & 6 & 6 & 0 & 6 & 0 & 0 & 0 \\
A & Meta-Harness & 1 & 15 & 15 & 0 & 15 & 0 & 0 & 0 \\
A & Meta-Harness & 2 & 9 & 9 & 0 & 9 & 0 & 0 & 0 \\
A & MIPROv2 & 0 & 3 & 3 & 0 & 3 & 0 & 0 & 0 \\
A & MIPROv2 & 1 & 4 & 4 & 0 & 4 & 0 & 0 & 0 \\
A & SkillOpt & 0 & 85 & 5 & 80 & 5 & 80 & 0 & 0 \\
A & SkillOpt & 1 & 9 & 9 & 0 & 9 & 0 & 0 & 0 \\
A & SkillOpt & 2 & 53 & 13 & 40 & 13 & 40 & 0 & 0 \\
A & Trace2Skill & 0 & 7 & 0 & 7 & 7 & 0 & 0 & 7 \\
A & Trace2Skill & 1 & 6 & 0 & 6 & 5 & 1 & 0 & 5 \\
A & Trace2Skill & 2 & 14 & 0 & 14 & 14 & 0 & 0 & 14 \\
A & WikiSkill & 0 & 3 & 0 & 3 & 3 & 0 & 0 & 3 \\
A & WikiSkill & 1 & 2 & 0 & 2 & 2 & 0 & 0 & 2 \\
A & WikiSkill & 2 & 3 & 0 & 3 & 3 & 0 & 0 & 3 \\
\midrule
B & CapEvolve AO & 0 & 2 & 2 & 0 & 2 & 0 & 0 & 0 \\
B & One-shot control & 0 & 7 & 7 & 0 & 7 & 0 & 0 & 0 \\
B & One-shot control & 1 & 1 & 1 & 0 & 1 & 0 & 0 & 0 \\
B & CapEvolve HC & 0 & 2 & 2 & 0 & 2 & 0 & 0 & 0 \\
B & CapEvolve HC & 1 & 1 & 1 & 0 & 1 & 0 & 0 & 0 \\
B & CapEvolve HC & 2 & 1 & 1 & 0 & 1 & 0 & 0 & 0 \\
\midrule
C & One-shot control & 1 & 1 & 1 & 0 & 1 & 0 & 0 & 0 \\
C & One-shot control & 2 & 1 & 1 & 0 & 1 & 0 & 0 & 0 \\
C & GEPA & 2 & 1 & 1 & 0 & 0 & 0 & 1 & 0 \\
\bottomrule
\end{tabular}
\end{table*}


\section{Cost and Token Totals}
\label{app:cost}

\begin{table}[h]
\caption{Metered usage so far (in progress; final totals follow the search cutoff). Tokens exclude cache reads. Calls include smoke tests, restarts and waste, which the metrics report separately. USD is a Claude-only list-price proxy and a lower bound.}
\label{tab:cost}
\centering\small
\begin{tabular}{@{}lrr@{}}
\toprule
Scope & Calls & Non-cache tokens (M) \\
\midrule
Suite A & \costACalls{} & \costATokensM{} \\
Suite B & \costBCalls{} & \costBTokensM{} \\
Suite C & \costCCalls{} & \costCTokensM{} \\
All, incl.\ smoke and untagged & \costCalls{} & \costTokensM{} \\
\bottomrule
\end{tabular}
\end{table}

Of the total, \costRitsM{}M tokens were served by RITS (uncapped by user decision) and \costNonRitsM{}M by capped gateway models. The proxy-priced total is \$\costUsdProxy{}. \needs{Final per-arm cost table: tokens by role, proxy USD, wall hours and cost per accepted candidate, from analysis/out/metrics after the cutoff.}
